\begingroup
\setlength{\parindent}{0pt}
\setlength{\parskip}{0.40\baselineskip}
\setlength{\abovedisplayskip}{6.5pt plus 1.5pt minus 1.5pt}
\setlength{\belowdisplayskip}{6.5pt plus 1.5pt minus 1.5pt}
\setlength{\jot}{5pt}
\clubpenalty=10000
\widowpenalty=10000

\section*{How the Proof Works}

We start with a smooth, incompressible flow of finite energy and keep its
viscosity fixed and positive. Its velocity \(u\) and pressure \(p\) satisfy
the three-dimensional Navier--Stokes equation:
\[
 \underbrace{\partial_tu}_{\text{local acceleration}}
 +\underbrace{(u\!\cdot\nabla)u}_{\text{nonlinear transport}}
 =\underbrace{-\nabla p}_{\text{pressure}}
 +\underbrace{\nu\Delta u}_{\text{viscous diffusion}},
 \qquad \nabla\!\cdot u=0.
\]
We know it stays smooth for a short time; the question is whether it stays
smooth forever or becomes singular in finite time. The hope is simple:
viscosity smooths the fluid faster than nonlinear transport can sharpen it.

Consider the case of a vortex within the fluid. Its vorticity
\(\omega=\nabla\times u\) satisfies
\[
 \partial_t\omega+(u\cdot\nabla)\omega
 =(\omega\cdot\nabla)u+\nu\Delta\omega.
\]
Strain within the fluid's flow can naturally stretch a vortex. When that
happens, incompressibility forces a material tube to narrow as it lengthens.
For a thin tube that remains locally circular, a positive axial strain
\(\sigma\) increases its length \(\ell\) and reduces its radius \(r\):
\[
 \dot\ell=\sigma\ell>0,
 \qquad r\propto\ell^{-1/2}.
\]
The fluid now circulates around a smaller radius. If the enclosed circulation
\(\Gamma\) is retained, its mean swirl speed \(v\) increases:
\[
 v=\frac{|\Gamma|}{2\pi r}.
\]

But the tube follows a chosen group of fluid particles, while the vortex
core is the region where rotation is concentrated. Stretching squeezes
those particles inward. Viscosity diffuses momentum between neighbouring
regions of fluid, reducing velocity differences and spreading the rotation
through the surrounding fluid. Over an elapsed time \(\tau\), its diffusion
distance is of order
\[
 r_{\mathrm{diff}}\sim\sqrt{\nu\tau}.
\]
The material tube can narrow while some rotation spreads beyond its boundary,
reducing the circulation enclosed by it. So narrowing makes the fluid
circulate faster only if the enclosed circulation does not weaken too
quickly. Whether viscosity can smooth those differences fast enough is the
difficulty.

If this concentration kept going, the region of fast motion would become
smaller and smaller. We would have to follow it toward vanishing width,
where even very large velocities could occupy very little volume. Looking
at the total energy would tell us less and less about what was happening
inside that region. To compare transport and diffusion there, we rescale
the flow: shrink lengths by \(\lambda>1\), increase velocities by
\(\lambda\), and shorten times by \(\lambda^2\):
\[
 \begin{aligned}
 u_\lambda(x,t)&=\lambda u(\lambda x,\lambda^2t),\\
 p_\lambda(x,t)&=\lambda^2p(\lambda x,\lambda^2t).
 \end{aligned}
\]
Every term in the momentum equation acquires a factor \(\lambda^3\).
Removing it leaves the original equation with the viscosity still fixed
at \(\nu\). A narrower, faster flow has preserved the scaling balance
between transport and diffusion. Shrinking the region has given viscosity
no automatic advantage; their actual competition still depends on how the
flow is arranged.

\needspace{8\baselineskip}
We need a measurement that can see the spatial variation inside this
shrinking region. The Sobolev index \(s\) determines how strongly we count
fine detail. At integer values it counts spatial derivatives in \(L^2\);
fractional values fill in between them by weighting frequency \(\xi\) by
\(|\xi|^s\). Under rescaling, velocity contributes \(\lambda\), the
derivative weight contributes \(\lambda^s\), and the square root of the
volume contributes \(\lambda^{-3/2}\). Thus
\[
 \|u_\lambda(\cdot,t)\|_{\dot H^s}
 =\lambda^{s-1/2}\|u(\cdot,\lambda^2t)\|_{\dot H^s}.
\]
Below \(s=\tfrac12\), concentration makes the norm smaller. Above
\(s=\tfrac12\), it makes the norm larger. Exactly at \(s=\tfrac12\),
the finer variation offsets the shrinking volume and the norm is unchanged.
This is the critical height: the measurement retains its sensitivity as
we change scale. That invariance compares rescaled flows; it does not say
the norm stays constant as one flow evolves.

Energy sits below that height. With \(E=\tfrac12\|u\|_2^2\),
\[
 E[u_\lambda](t)=\lambda^{-1}E[u](\lambda^2t).
\]
The velocity grows, yet the volume containing it shrinks fast enough for
the total energy to fall. An energy bound can coexist with increasingly
fast motion in an increasingly narrow region. It cannot, by itself, stop
this concentration.

The time needed for each further stage also gets shorter. At width \(r\)
and velocity scale \(U\), transport acts over a time of order \(r/U\),
while viscosity acts over a time of order \(r^2/\nu\). Both are shortened
by \(\lambda^2\) under the rescaling. If a flow could repeat such stages,
their durations would fit into a finite total:
\[
 \tau_j=\lambda^{-2j}\tau_0,
 \qquad
 \sum_{j=0}^{\infty}\tau_j
 =\frac{\tau_0}{1-\lambda^{-2}}<\infty.
\]
This is the hypothetical Zeno cascade. Infinitely many narrowing stages
could accumulate before a finite time \(T_*\), with no smallest width
and no final stage of acceleration. The scaling permits that picture;
whether the fluid can actually carry it out is the regularity problem.

For a finite-time singularity, velocity would have to blow up. That requires
acceleration to blow up, then jerk, and so on, and so on. To see the necessity,
follow the velocity field at fixed spatial positions. Its temporal responses
are \(u,\partial_tu,\partial_t^2u,\ldots\), and for any \(t_0<T_*\),
\[
 \|\partial_t^ju(t)\|_\infty
 \le \|\partial_t^ju(t_0)\|_\infty
 +\int_{t_0}^{t}\|\partial_t^{j+1}u(s)\|_\infty\,ds.
\]
A bounded next derivative could produce only a finite change in the time
remaining. Unbounded velocity requires an infinite accumulated temporal
response, which forces that response to become unbounded too. Apply this
again and the requirement reaches every finite temporal order. These are
Eulerian responses; the acceleration of a moving fluid particle is
\(D_tu=\partial_tu+(u\cdot\nabla)u\). The claim concerns unbounded spatial
norms on the approach to \(T_*\), without requiring monotone growth or
growth at one fixed point.

\needspace{11\baselineskip}
Those increasingly demanding temporal responses are generated by the
Navier--Stokes equation itself. Differentiating it lets us see what the
fluid would have to supply at each order:
\[
 \begin{aligned}
 \partial_tu={}&\nu\Delta u-(u\cdot\nabla)u-\nabla p,\\
 \partial_t^2u={}&\nu\Delta\partial_tu
 -(\partial_tu\cdot\nabla)u\\
 &-(u\cdot\nabla)\partial_tu-\nabla\partial_tp,\\
 \partial_t^3u={}&\nu\Delta\partial_t^2u
 -(\partial_t^2u\cdot\nabla)u\\
 &-2(\partial_tu\cdot\nabla)\partial_tu
 -(u\cdot\nabla)\partial_t^2u
 -\nabla\partial_t^2p.
 \end{aligned}
\]
Each differentiation splits the transport product again, while viscosity
applies a Laplacian to the preceding response. The second equation contains
\(\Delta\partial_tu\); substituting the first equation into it exposes
\(\Delta^2u\). The next substitution exposes \(\Delta^3u\). Higher time
derivatives keep reaching further into the spatial structure of the flow.
The complete pattern retains all the accompanying nonlinear terms:
\[
 \begin{aligned}
 \partial_t^{n+1}u={}&\nu\Delta\partial_t^nu\\
 &-\sum_{a=0}^{n}\binom na
 (\partial_t^au\cdot\nabla)\partial_t^{n-a}u
 -\nabla\partial_t^np.
 \end{aligned}
\]
We write these velocity and pressure responses as
\(V_n=\partial_t^nu\) and \(Q_n=\partial_t^np\).

At the critical height, this temporal-to-spatial relation becomes a relation
between Sobolev energies. Write
\(E_s=\tfrac12\|\Lambda^su\|_2^2\), where
\(\Lambda=(-\Delta)^{1/2}\), and let \(H=E_{1/2}\).
Differentiating \(E_s\) pairs the velocity with its temporal response.
The viscous Laplacian moves one derivative onto each factor under integration
by parts, exposing the next spatial height:
\[
 E_s'=-2\nu E_{s+1}+\mathcal P_s,
\]
where the transport and pressure contribution is retained as
\[
 \mathcal P_s=-\operatorname{Re}\langle\Lambda^su,
 \Lambda^s((u\cdot\nabla)u+\nabla p)\rangle_{L^2}.
\]
Differentiate \(H\) once and \(E_{3/2}\) appears. Differentiate again
and the changing \(E_{3/2}\) exposes \(E_{5/2}\):
\[
 \begin{aligned}
 H'&=-2\nu E_{3/2}+\mathcal P_{1/2},\\
 H''&=4\nu^2E_{5/2}-2\nu\mathcal P_{3/2}
       +\partial_t\mathcal P_{1/2}.
 \end{aligned}
\]
Continuing through \(H''',H^{(4)},\ldots\) exposes
\(E_{7/2},E_{9/2},\ldots\), together with the differentiated nonlinear
contributions. The temporal sequence that a singularity would force us
to confront has led directly to an unending sequence of spatial Sobolev
heights.

And those heights are exactly what smoothness needs. For any chosen number
\(k\) of continuous spatial derivatives, Sobolev embedding asks for a
finite height above \(k+\tfrac32\):
\[
 m>k+\tfrac32
 \quad\Longrightarrow\quad H^m(\mathbb R^3)\hookrightarrow C^k(\mathbb R^3).
\]
If every finite height is available, there is no last derivative we can
ask for. The intersection contains enough regularity for every choice of
\(k\):
\[
 L^2\cap\bigcap_{n\ge0}\dot H^{n+1/2}
 =\bigcap_{m\ge0}H^m
 =H^\infty\hookrightarrow C^\infty.
\]
\needspace{8\baselineskip}
This is where the temporal and spatial demands meet. The equation connects
each new temporal response to higher spatial derivatives, and the
intersection of all those spatial requirements is smoothness. With those
requirements controlled together up to \(T_*\), the limiting velocity
would be smooth enough to start the flow again. The identities expose
that possibility; realizing it requires bounds on the full approach to
\(T_*\), including the nonlinear contributions and agreement between
all the retained derivatives.

If a singularity were to form, we would expect the opposite behavior.
The flow could remain smooth at every earlier instant while some fixed
spatial Sobolev norm became unbounded on the approach to \(T_*\).
Its temporal responses would become unbounded too, even while total
energy remained finite. Every shorter interval could contain all finite
derivatives, yet their control would deteriorate as its endpoint approached
the singular time. We would lose the regularity needed to carry the flow
through \(T_*\). That is the loss the construction has to exclude.

We begin that construction with the initial velocity
\(u_0\in H^\infty_\sigma(\mathbb R^3)\): every finite spatial Sobolev
order is available at the start, and the velocity is incompressible.
Incompressibility fixes pressure through
\(-\Delta p=\partial_i\partial_j(u_i u_j)\), with its decay normalization.
Differentiating this relation along with the velocity determines the
pressure response at each order. In the notation already introduced,
\[
 \begin{aligned}
 V_{n+1}&=\nu\Delta V_n
 -\sum_{a=0}^n\binom na(V_a\cdot\nabla)V_{n-a}-\nabla Q_n,\\
 Q_n&=R_iR_j\sum_{a=0}^n\binom na V_{a,i}V_{n-a,j}.
 \end{aligned}
\]
Here \(R_i\) are the Riesz transforms and repeated spatial indices are
summed. At \(t=0\), \(V_0=u_0\) fixes the first response, which fixes the
next, and so on. Spatial differentiation fills out each temporal row.
These velocity and pressure derivatives form the mixed jet generated by
the initial datum. Its entries remain connected through the equation
that generated them.

Any finite regularity demand uses only a finite part of this structure.
A temporal depth \(N\) and a spatial height \(M\) select a rectangle
\(\mathcal R_{N,M}[u_0;T_*]\). Each coordinate measures a temporal response
\(V_n\) in a spatial Sobolev space; adjacent temporal rows retain
\(\partial_tV_n=V_{n+1}\), and the pressure rows accompany them.
The datum-generated rectangle theorem, Theorem~\ref{thm:datum-rectangles},
controls this finite selection throughout the approach to \(T_*\):
\[
 \sup_{0<T<T_*}\sum_{n=0}^{N}\left(
 \|V_n\|_{L^2(0,T;H^{M+1})}^2
 +\|V_{n+1}\|_{L^2(0,T;H^{M-1})}^2
 \right)<\infty.
\]
For fixed \(N,M\), the bound stays finite as \(T\) approaches \(T_*\).
It may increase when we ask for more derivatives. What it excludes is
divergence at any fixed pair of orders as the proposed singular time
approaches.

The adjacent spaces let each retained response reach that time.
The upper row \(V_n\in L^2(0,T_*;H^{M+1})\) and its temporal derivative
\(V_{n+1}\in L^2(0,T_*;H^{M-1})\) pair across \(H^M\).
Their square-integrability makes the change in the squared \(H^M\) norm
integrable. The Hilbert-triple trace theorem then gives
\[
 V_n\in C([0,T_*];H^M).
\]
So each finite rectangle supplies actual endpoint values at its retained
orders. A larger rectangle keeps every derivative already present, and
on an overlap both rectangles contain derivatives of the original velocity
and pressure. Their entries agree, including their traces at \(T_*\).

\needspace{7\baselineskip}
This agreement lets us retain every finite rectangle together. Increasing
temporal depth and spatial height independently gives the compatible
intersection
\[
 \begin{aligned}
 \mathcal I[u_0]
 &:=\bigl(\mathcal R_{N,M}[u_0;T_*]\bigr)_{N,M\ge0},\\
 u&\in\bigcap_{r,m\ge0}H^r(0,T_*;H^m(\mathbb R^3)).
 \end{aligned}
\]
Every finite request for derivatives fits inside a sufficiently large
rectangle, and every larger rectangle preserves the answer. We do not
need one bound that works uniformly over all orders. Each finite order
has its own bound, and compatibility makes them properties of one flow.

\needspace{8\baselineskip}
In particular, the zeroth temporal row reaches \(T_*\) in every spatial
Sobolev space. Its traces agree as a single limiting velocity:
\[
 \begin{aligned}
 u_*&=\lim_{t\uparrow T_*}u(t)\quad\text{in every }H^m,\\
 u_*&\in H^\infty_\sigma\hookrightarrow C^\infty.
 \end{aligned}
\]
The subscript \(\sigma\) records incompressibility. The differentiated
equations also reach the endpoint, fixing its temporal responses from
\(u_*\). Smooth local existence starts a solution there, and uniqueness
joins it to the original flow. The flow continues beyond its supposed
maximal time. A finite \(T_*\) is impossible, so \(T_*=\infty\).

With the intersection constructed, we can derive the finite estimates,
and use them. For any finite time \(T\), temporal order \(r\), spatial
order \(k\), and \(2\le p\le\infty\), choose
\(m>k+\tfrac32-\tfrac3p\). A sufficiently large rectangle and Sobolev
embedding give
\[
 \partial_t^ru\in C([0,T];H^m)
 \hookrightarrow L^\infty(0,T;W^{k,p}).
\]
For example, \(r=0\), \(k=1\), \(p=\infty\), and \(m=3\) give
\[
 \int_0^T\|\nabla u(t)\|_\infty\,dt
 \le T\sup_{0\le t\le T}\|\nabla u(t)\|_\infty<\infty.
\]
Return to the locally circular material tube at the start.
Its axial strain satisfies \(\sigma\le\|\nabla u\|_\infty\), while
incompressibility gives \(\dot r/r=-\sigma/2\). Integrating,
\[
 r(t)\ge r(0)\exp\!\left(
 -\frac12\int_0^t\|\nabla u(s)\|_\infty\,ds
 \right)>0.
\]
Only a finite amount of stretching can accumulate in a finite time.
The tube cannot be squeezed to zero radius there.

\endgroup
